\documentclass[12pt,reqno]{amsart}
\usepackage{math-review}
% Project-wide mathematical notation. No claims are certified here.

\title{The Schur Complement -- Lecture}
\author{Known-mathematics exposition benchmark}
\date{}
\begin{document}
\maketitle
\section{Goal theorem and the strategy}
We will derive the block positivity test from ordinary completion of a
square. The unfamiliar object will appear as a remainder, not as a formula
to memorize in advance.
\begin{theorem}[Classical Schur-complement criterion]\label{thm:main}
Let $p,q\geq1$ be integers. Let $A\in\mathbb R^{p\times p}$ and
$C\in\mathbb R^{q\times q}$ be symmetric, and let $B\in\mathbb R^{p\times q}$.
Assume that $x^{\mathsf T}Ax>0$ for every $0\ne x\in\mathbb R^p$. Define
\[
 M=\begin{pmatrix}A&B\\ B^{\mathsf T}&C\end{pmatrix},
 \qquad S=C-B^{\mathsf T}A^{-1}B.
\]
Then $M$ is positive definite if and only if $S$ is positive definite.
Here a real symmetric matrix $H$ is positive definite when
$v^{\mathsf T}Hv>0$ for every nonzero real vector $v$ of the appropriate size.
\end{theorem}
The proof will remove the mixed term, check that nonzero vectors are
preserved, and test the residual form in both directions.

\section{A scalar model motivates the shift}
Assume the reader can multiply matrices and test positivity by a quadratic
form. Begin with $a>0$ and the scalar expression
\[
 ax^2+2bxy+cy^2
 =a\left(x+\frac ba y\right)^2+\left(c-\frac{b^2}{a}\right)y^2.
\]
The coefficient $b/a$ is forced by matching the cross term. In blocks,
the matching equation is $AK=B$, so the shift is $K=A^{-1}B$.

\section{Constructing and checking the identity}
Write $z=x+A^{-1}By$. Since $A$ is symmetric and invertible,
$(A^{-1})^{\mathsf T}=A^{-1}$. Expanding every term gives
\begin{align*}
 z^{\mathsf T}Az
 &=x^{\mathsf T}Ax+x^{\mathsf T}By+y^{\mathsf T}B^{\mathsf T}x
       +y^{\mathsf T}B^{\mathsf T}A^{-1}By\\
 &=x^{\mathsf T}Ax+2x^{\mathsf T}By
       +y^{\mathsf T}B^{\mathsf T}A^{-1}By.
\end{align*}
The two middle terms agree because they are transposes of the same real
scalar. Moving the last term to the remainder proves the precise identity.
\begin{lemma}\label{lem:square}
Let $A\in\mathbb R^{p\times p}$ be symmetric and invertible,
$B\in\mathbb R^{p\times q}$ and $C\in\mathbb R^{q\times q}$ symmetric.
For all $x\in\mathbb R^p$ and $y\in\mathbb R^q$, putting
$z=x+A^{-1}By$ gives
\[
 x^{\mathsf T}Ax+2x^{\mathsf T}By+y^{\mathsf T}Cy
 =z^{\mathsf T}Az+y^{\mathsf T}(C-B^{\mathsf T}A^{-1}B)y.
\]
\end{lemma}
Denote the left side of this identity by $Q(x,y)$.

The identity was derived with invertibility; the sign of $A$ is not used
until we turn it into a positivity argument.

\section{Preserving vectors and closing the proof}
We must still justify that the new variables test all nonzero pairs.
\begin{lemma}\label{lem:inverse}
For every real matrix $K\in\mathbb R^{p\times q}$, the linear map
\[
 (x,y)\longmapsto (x+Ky,y)
\]
is invertible, with inverse $(z,y)\mapsto(z-Ky,y)$. Consequently it takes
nonzero vectors to nonzero vectors, in both directions.
\end{lemma}
To prove this, substitution of $x=z-Ky$ into $x+Ky$ gives $z$;
substitution of $z=x+Ky$ into $z-Ky$ gives $x$. In particular,
$(z,y)=(0,0)$ forces $(x,y)=(0,0)$, and conversely.

\begin{proof}[Proof of Theorem~\ref{thm:main}]
For sufficiency, assume $S>0$ in addition to the given $A>0$.
For any nonzero $(x,y)$ the new pair $(z,y)$ is nonzero, and therefore
\[
 Q(x,y)=z^{\mathsf T}Az+y^{\mathsf T}Sy>0.
\]
If $z\ne0$, the first term is positive and the other is nonnegative;
if $z=0$, then $y\ne0$ and the second term is positive.

For necessity, assume $M>0$ and let $y$ be any nonzero vector. Set
$x=-A^{-1}By$. The pair $(x,y)$ is nonzero, so
\[
 y^{\mathsf T}Sy=Q(-A^{-1}By,y)>0.
\]
This universal test is precisely the positive definiteness of $S$.
\end{proof}
No external block-matrix theorem was imported into this proof. The scalar
model motivated the construction; the algebra and the inverse map closed it.

\section{Using the criterion and testing its boundary}
For $A=2$, $B=1$, $C=1$ the remainder is $S=1/2$. Indeed,
\[
 2x^2+2xy+y^2=2(x+y/2)^2+y^2/2>0
 \quad\text{for }(x,y)\ne(0,0).
\]
For fixed $y$, this expression has minimum $y^2/2$ at $x=-y/2$.
This computes the elimination, rather than merely invoking its name.

Now take $A=-1$, $B=0$, $C=1$. The identity still holds and $S=1>0$,
but the value at $(x,y)=(1,0)$ is $-1$. Thus the step using the sign of
the first square fails exactly when the positive-block assumption is removed.
If $A$ is singular, even the shift is undefined; that case is not covered.

\nocite{S-guide}
\bibliographystyle{plain}
\bibliography{../generated/references}
\end{document}
